\par\addvspace{12pt}\Needspace{9\baselineskip}

\section{随机变量}

% ITEM: probability-example-014
\solutionitem{probability-example-014}{例6.14}
独立抛掷两枚公平骰子，\(X\)为较大点数。求\(X\)与\(6-X\)的分布列。
\solutionanswer
\(P(X\le k)=k^2/36\)，所以\(P(X=k)=[k^2-(k-1)^2]/36=(2k-1)/36\)，\(k=1,\ldots,6\)。因此\(6-X\)在\(0,1,2,3,4,5\)处的概率依次为\(11,9,7,5,3,1\)各除以\(36\)。变换后的数值次序与原变量相反。


% ITEM: probability-example-015
\solutionitem{probability-example-015}{例6.15}
本题所指分布函数为
\[
 F_Y(y)=\begin{cases}0,&y<0,\\1/4,&0\le y<1,\\3/4,&1\le y<2,\\1,&y\ge2.\end{cases}
\]

利用上述\(F_Y\)，求\(P(Y<1)\)、\(P(Y\le1)\)与\(P(1\le Y\le2)\)，再从各跳跃高度恢复\(Y\)的分布列。
\solutionanswer
\(P(Y<1)=F_Y(1^-)=1/4\)，\(P(Y\le1)=F_Y(1)=3/4\)，\(P(1\le Y\le2)=F_Y(2)-F_Y(1^-)=3/4\)。在\(0,1,2\)处的跳跃高度分别为\(1/4,1/2,1/4\)，即各点概率。严格与非严格端点在离散分布中不能互换。


% ITEM: probability-exercise-08-01
\solutionitem{probability-exercise-08-01}{习题6.8（1）}
随机变量\(X\)在\(-1,0,1\)处的概率分别为\(1/4,1/2,1/4\)。写出\(X^2\)和\(2X+1\)的分布列及\(X\)的分布函数。
\solutionanswer
\(X^2\)取\(0,1\)的概率均为\(1/2\)；\(2X+1\)取\(-1,1,3\)的概率依次为\(1/4,1/2,1/4\)。分布函数为
\[
 F_X(x)=\begin{cases}0,&x<-1,\\1/4,&-1\le x<0,\\3/4,&0\le x<1,\\1,&x\ge1.\end{cases}
\]


% ITEM: probability-exercise-08-02
\solutionitem{probability-exercise-08-02}{习题6.8（2）}
从编号1至4的球中不放回抽两个，令\(X\)为第一次的数字，\(Y\)为第二次的数字。写出联合分布，求\(P(X+Y=5)\)，并判断独立性。
\solutionanswer
联合概率\(P(X=i,Y=j)=1/12\)对\(1\le i,j\le4,i\ne j\)，对\(i=j\)为零。和为5对应\((1,4),(2,3),(3,2),(4,1)\)，概率\(1/3\)。边际各均匀，但\(P(X=1,Y=1)=0\ne1/16\)，故不独立。


% ITEM: probability-exercise-08-03
\solutionitem{probability-exercise-08-03}{习题6.8（3）}
本题独立公平抛两次硬币，\(X\)为第一次正面的指示量，\(Y\)为总正面次数；\((X,Y)=(0,0),(0,1),(1,1),(1,2)\)各以\(1/4\)发生。
在本节抛币联合表中，求\(X+Y\)的分布及\(P(X=1\mid Y\ge1)\)。
\solutionanswer
四个抛币序列\(TT,TH,HT,HH\)对应\((X,Y)=(0,0),(0,1),(1,1),(1,2)\)，各概率\(1/4\)。所以\(X+Y\)在\(0,1,2,3\)处的概率均为\(1/4\)。条件\(Y\ge1\)保留后三种，其中\(X=1\)有两种，故条件概率\(2/3\)。


% ITEM: probability-exercise-08-04
\solutionitem{probability-exercise-08-04}{习题6.8（4）}
设\(X,Y\)独立，各以\(1/2\)的概率取\(-1,1\)。求\(X+Y\)、\(XY\)的分布，判断\(X\)与\(XY\)是否独立。
\solutionanswer
和\(X+Y\)在\(-2,0,2\)处的概率为\(1/4,1/2,1/4\)，积\(XY\)在\(-1,1\)处的概率均为\(1/2\)。对每一组\((X,XY)=(\pm1,\pm1)\)，恰有一组原取值\((X,Y)\)与之对应，联合概率均为\(1/4\)，等于边际乘积，所以\(X\)与\(XY\)独立。


% ITEM: probability-exercise-08-05
\solutionitem{probability-exercise-08-05}{习题6.8（5）}
（2025年全国II卷第19题，题面整理）甲、乙进行乒乓球练习，每球胜者得1分，负者得0分。甲每球获胜的概率为\(p\)，乙为\(q=1-p\)，其中\(1/2<p<1\)，各球胜负相互独立。记\(p_k,q_k\)分别为打完\(k\ge2\)球后，甲、乙比对方至少多得2分的概率。

求\(p_3,p_4\)，并在\((p_4-p_3)/(q_4-q_3)=4\)时求\(p\)。再证明，对任意正整数\(m\)，
\[
 p_{2m+1}-q_{2m+1}<p_{2m}-q_{2m}<p_{2m+2}-q_{2m+2}.
\]
\solutionanswer
三球后至少领先2分只能是三球全胜，故\(p_3=p^3\)；四球后须至少胜三球，故\(p_4=p^4+4p^3q\)。同理\(q_3=q^3,q_4=q^4+4q^3p\)，因此
\[
 \frac{p_4-p_3}{q_4-q_3}=\frac{3p^3q}{3q^3p}=\frac{p^2}{q^2}=4,
\]
得\(p=2/3\)。这些概率也可先按胜球位置直接分类计数，无须预先使用分布名称。

令\(S_n\)为甲减乙的分差，\(D_n=p_n-q_n\)。从偶数\(2m\)步增加一球，只有分差\(2,-2\)可能失去原有的至少2分优势，故
\[
 D_{2m}-D_{2m+1}=qP(S_{2m}=2)-pP(S_{2m}=-2)
 =\binom{2m}{m+1}(pq)^m(p-q)>0.
\]
从奇数步增加一球，则只有分差\(1,-1\)可能新进入优势事件，故
\[
 D_{2m+2}-D_{2m+1}=pP(S_{2m+1}=1)-qP(S_{2m+1}=-1)
 =\binom{2m+1}{m}(pq)^m(p-q).
\]
将后一式减去前一式，并用Pascal关系与组合数对称性，得\(D_{2m+2}-D_{2m}=\binom{2m}{m}(pq)^m(p-q)>0\)。两项严格比较合起来正是所求不等式。
\par
